% Adaptador único del capítulo 38 del sucesor 102.
% Sustituye exclusivamente la jerarquía de encabezados. Los cuerpos de los
% dos propietarios científicos y del delta Ley 9 se conservan línea a línea,
% reordenados conforme al DAG causal aprobado. Véase MAPA_SEGMENTOS_CAPITULO_38.tsv.

% HMT_SOURCE_SEGMENT_BEGIN id=B00 source=colaboracion/parte_iii/source/public_final/base_83/c30_body.tex body_lines=1-2
\label{chap:p3-83-30-angulo_contraangulo}

% HMT_SOURCE_SEGMENT_END id=B00
% HMT_SOURCE_SEGMENT_BEGIN id=A00 source=colaboracion/parte_iii/source/public_final/ascensos_20260826/16f_realizacion_analitica_contraangulo.tex body_lines=1-4
\label{p3-20260826:chap:realizacion-analitica-contraangulo}
\index{counter-angle}
\index{microfiber of pi@microfiber of \(\pi\)!analytic character}
\index{determinantal character}
% HMT_SOURCE_SEGMENT_END id=A00
% HMT_SOURCE_SEGMENT_BEGIN id=D00 source=manuscrito/sucesor_102/deltas_ley9/c38_par_angular_cartas_y_canales.tex body_lines=1
% Distribución causal exacta del delta Ley 9; contenido conservado sin síntesis.
% HMT_SOURCE_SEGMENT_END id=D00

\section{Intrinsic angular phase \texorpdfstring{\(\alpha_{\mathrm{HMT}}\)}{alpha HMT}, nonadic lift \texorpdfstring{\(10^3\)}{10^3}, and precursor \texorpdfstring{\(x_A\)}{x_A}}
\label{sec:s102:c38:1}

\subsection{Construction of the analytic precursor \texorpdfstring{\(x_A\)}{x_A}, the contraction \texorpdfstring{\(D_A\)}{D_A}, and the regional correction \texorpdfstring{\(\mathscr C_\pi(\alpha)\)}{C_pi(alpha)}}
\label{sec:s102:c38:1:1}

% HMT_SOURCE_SEGMENT_BEGIN id=B01 source=colaboracion/parte_iii/source/public_final/base_83/c30_body.tex body_lines=4-32
\label{p3-83-c30-s1:chap:realizacion-analitica-contraangulo}
\index{conjugate angular coordinate \(C^*\)}
\index{microfiber of pi@microfiber of \(\pi\)!analytic character}
\index{determinantal character}

The finite catalog and an analytic evaluation answer questions of different
types. The catalog determines which regions exist, which share a
return and what information each preserves. An analytic evaluation must
also decide how the length of a word is transformed into degree, how
a branch is prolonged after its return and with which scalar character that
prolongation is read. This chapter explicitly constructs that second
operation.

The construction begins with four data already obtained in the preceding
chapters: the five-region microfiber of \(\pi\), the length six of
each regional word, the dodecaphasic closure of \(\alpha\), and the bilayer
transport associated with the common coordinate
\(A=10^3\alpha\). No value of the conjugate angular coordinate \(C_*^{\rm HMT}\), of the reduced action
component, of the Barbero--Immirzi functional, or of a metrological constant is
used in the evaluation. The result is a convergent regional character,
an exact recurrence, and an oriented phase whose sexagesimal chart reproduces
the published conjugate angular coordinate.

The separation between the finite part and the analytic part is essential. The
integer \(169\) will be a theorem of the catalogue. The series that prolongs it will be the
theorem of an analytic reader whose rules are declared before it is evaluated.
Thus, the additional premises are not concealed within a decimal
coincidence.

% HMT_SOURCE_SEGMENT_END id=B01
\section{Preangular construction of \texorpdfstring{\(\eta_{\mathrm{ret}}\)}{eta_ret} through \texorpdfstring{\(H_5\)}{H5}, \texorpdfstring{\(D_A\)}{D_A} and \texorpdfstring{\(\mathscr C_\pi(\alpha)\)}{C_pi(alpha)}}
\label{sec:s102:c38:2}

\subsection{Independent derivation of \texorpdfstring{\(\eta_{\mathrm{ret}}=H_5-(90/\pi)\mathscr C_\pi(\alpha)\)}{eta_ret=H5-(90/pi)C_pi(alpha)} through the regional analytic reader}
\label{sec:s102:c38:2:1}

% HMT_SOURCE_SEGMENT_BEGIN id=A01 source=colaboracion/parte_iii/source/public_final/ascensos_20260826/16f_realizacion_analitica_contraangulo.tex body_lines=5-38,40-47

The finite catalogue and an analytic evaluation answer questions of different types. The catalogue determines which regions exist, which share a return, and what information each preserves. An analytic evaluation must further determine how the length of a word is transformed into degree, how a branch is prolonged after its return, and by which scalar character that prolongation is read. This section constructs that second operation explicitly.

The construction begins with the five-region microfiber of \(\pi\), the
length six of each regional word, the dodecaphasic closure of \(\alpha\),
and the declared angular insertion
\[
 \iota_\angle(x):=10^3x\,u_\circ,\qquad
 \Adeg:=10^3\alpha .
\]
Here \(u_\circ\) is the sexagesimal base whose complete turn has
coordinate \(360\).
The factor \(10^3\) belongs to that chart normalization and is not deduced from the
arithmetic closure of \(\alpha\). The two-sheet transport is subsequently constructed
from \(\Adeg\). The function evaluated by the reader does not receive as
an argument the counter-angle, the post-return coordinate, the
hexad--octad functional, or a metrological constant. The result is a
convergent regional character, an exact recurrence, and an oriented phase in the
declared sexagesimal chart.

The separation between the finite part and the analytic part is essential. The
integer \(169\) is a theorem of the catalogue. The series that prolongs it belongs
to an analytic reader whose rules are declared before evaluating it.

The phase is denoted by \(y_*\) and, when its regional origin is emphasized, by
\(y_{\rm reg}\); both symbols denote the same reader constructed
below.


The evaluation receives the catalogue of \(468\) emissions, the exact value of
\(\alpha_{\rm HMT}\), \(\varphi=(1+\sqrt5)/2\), the fixed formula of \(H_5\),
the angular chart, and the five rules of the regional reader. With those data, it
computes \(169\), \(D_A\), the recurrence, and \(C^*\). The theorem is exact in
this declared class of readers; it does not use \(C^*\) as an input or assert that
the rules of \(H_5\) are the only possible natural class.

% HMT_SOURCE_SEGMENT_END id=A01
\subsection{Persistence of the return in the \texorpdfstring{\(5\)}{5} regions of the microfiber}
\label{sec:s102:c38:2:2}

% HMT_SOURCE_SEGMENT_BEGIN id=B02 source=colaboracion/parte_iii/source/public_final/base_83/c30_body.tex body_lines=34-155

Let
\[
\mathcal F_\pi=\Psi^{-1}(010211)\subset\mathcal U_6
\]
the microfiber relative to the catalog. Recall its five elements, now
separated into initial block and terminal block:
\[
\begin{aligned}
 &(501,614,498\mid169,272,272),\\
 &(810,923,870\mid169,272,272),\\
 &(810,983,810\mid169,272,272),\\
 &(870,923,810\mid169,272,272),\\
 &(870,923,810\mid418,674,521).
\end{aligned}
\]
We define the return projection.
\[
p_{\rm ret}:\mathbb Z^6\longrightarrow\mathbb Z^3,
\qquad
p_{\rm ret}(u_1,\ldots,u_6)=(u_4,u_5,u_6).
\]

\begin{proposition}[Persistent terminal block]
\label{p3-83-c30-s1:prop:bloque-terminal-persistente}
The multiset \(p_{\rm ret}(\mathcal F_\pi)\) has a unique element of
maximum multiplicity:
\[
b_\pi=(169,272,272),
\qquad
\operatorname{mult}_{\mathcal F_\pi}(b_\pi)=4.
\]
The only remaining terminal block is
\[
b_\pi^-=(418,674,521)
\]
and has multiplicity one.
\end{proposition}

\begin{proof}
The assertion follows by exact enumeration of the five rows of the
catalog that satisfy \(\Psi(U)=010211\). Four rows have terminal
block \((169,272,272)\); the fifth, which carries the mutated return, has
\((418,674,521)\). The maximum multiplicity is four and is attained only
once.
\end{proof}

The oriented difference between the mutated return and the persistent is
\[
 \omega_\pi=b_\pi^- - b_\pi=(249,402,249),
 \qquad
 \langle(1,1,1),\omega_\pi\rangle=900.
\]
This identity separates two objects that must not be conflated. The integer
\(169\) is a coordinate of a terminal datum, that is, a datum of degree
zero. The vector \(\omega_\pi\) can be read as a one-cochain
after orienting the edge joining the two blocks. Calling it a cocycle would additionally
require fixing a cochain complex and verifying that its coboundary vanishes.
The analytic prolongation constructed below does not presuppose that
identification.

The selection of \(169\) does not require privileging the first position of the block. The symmetric group \(\mathfrak S_3\) acts by permuting its three coordinates. The stabilizer of \(b_\pi\) exchanges the two occurrences of \(272\) and fixes the unique coordinate of multiplicity one. Denote that coordinate by \(\operatorname{sing}(b)\) when the block has type \((r,s,s)\), with \(r\ne s\). Equivalently,
\[
\operatorname{sing}(b)
=\sum_{j=1}^3b_j-2\operatorname{moda}(b).
\]
In particular,
\[
\boxed{\operatorname{sing}(b_\pi)=169.}
\]

\begin{definition}[Residual fiber selector]
\label{p3-83-c30-s1:def:selector-residual-pi}
The selector \(\rho_\pi\) successively applies two operations:
\begin{enumerate}
  \item selects the terminal block of maximum multiplicity within
  \(p_{\rm ret}(\mathcal F_\pi)\);
  \item retains the coordinate fixed by the stabilizer of that block.
\end{enumerate}
By \cref{p3-83-c30-s1:prop:bloque-terminal-persistente},
\[
\boxed{\rho_\pi=169.}
\]
\end{definition}

\begin{theorem}[Equivariant uniqueness of return]
\label{p3-83-c30-s1:thm:unicidad-retorno-169}
Among the selectors that
\begin{enumerate}
  \item are invariant under permutations of the five regions;
  \item are equivariant under permutations of the three terminal
  coordinates;
  \item select the block of maximal multiplicity;
  \item retain a coordinate fixed by its stabilizer,
\end{enumerate}
the return value is uniquely determined and is \(169\).
\end{theorem}

\begin{proof}
Invariance under \(\mathfrak S_5\) prevents using the order of the rows. The
third condition selects \(b_\pi\), which is the unique mode by
\cref{p3-83-c30-s1:prop:bloque-terminal-persistente}. Its stabilizer in
\(\mathfrak S_3\) has two orbits on the coordinates: the pair of
values \(272\) and the singleton value \(169\). The fourth condition selects the
second orbit. Equivariance preserves its value when its position is shifted.
\end{proof}

This argument improves the positional reading of the return. The \(169\) is not
«the first coordinate» of a selected row: it is an invariant of the
regional multiset and of the action of its reordering symmetries.
Nor is it selected by global rarity. Among the \(468\) emissions in the
catalogue, \(169\) appears \(84\) times, uniformly distributed with
\(14\) occurrences in each of the six positions. The uniqueness of
\cref{p3-83-c30-s1:thm:unicidad-retorno-169} is therefore relational and regional: it resides
in the combination of microfiber, modal persistence, and stabilizer, not in the
isolated frequency of the integer.

% HMT_SOURCE_SEGMENT_END id=B02
\subsection{Construction of the regional analytic reader and its convergence domain}
\label{sec:s102:c38:2:3}

% HMT_SOURCE_SEGMENT_BEGIN id=B05 source=colaboracion/parte_iii/source/public_final/base_83/c30_body.tex body_lines=261-301

Analytic continuation is fixed by the following rules.

\begin{definition}[Regional Determinantal Reader]
\label{p3-83-c30-s1:def:lector-regional-determinantal}
The determinantal regional reader satisfies:
\begin{enumerate}
  \item \emph{degree compatibility}: cylindrical length is the analytic
  degree;
  \item \emph{renewal decomposition}: the finite return impulse and
  strict continuation belong to distinct additive summands;
  \item \emph{continuation normalization}: after recording the residue
  \(\rho_\pi\) just once, the surviving branch begins at the unit of
  the determinant line;
  \item \emph{concatenation covariance}: each elementary prolongation
  acts through \(\bigwedge^2\mathcal T=D_A\);
  \item \emph{orientation}: in the cardinal convention fixed for APP, the
  regional return belongs to the compensating sector and is subtracted from the fifth-degree
action phase.
\end{enumerate}
\end{definition}

The third rule has mathematical content. The value \(169\) is recorded once as the amplitude of the return; it is not multiplied again at each prolongation. The strict continuation counts the sole line that continues and normalizes it by its unit. This is the distinction between a renewal recurrence and homogeneous iteration of the initial impulse.

The need to declare that normalization can be proved. For every
\(\lambda\in\mathbb R\), the family
\[
F_\lambda(z)
=169z^6+\lambda\frac{D_Az^7}{1-D_Az}
\]
preserves the same finite return and the same determinantal ratio between
consecutive coefficients. The catalogue and the ratio alone do not distinguish
\(\lambda\). The unit of the determinant line fixes
\(\lambda=1\) before \(z=\alpha\) is evaluated. Thus, the normalization
is not obtained by fitting the final value; it forms part of the definition of the
character.

% HMT_SOURCE_SEGMENT_END id=B05
\subsection{Analytical solution of the return recurrence}
\label{sec:s102:c38:2:4}

% HMT_SOURCE_SEGMENT_BEGIN id=B06 source=colaboracion/parte_iii/source/public_final/base_83/c30_body.tex body_lines=303-364

We define the coefficients \(\kappa_n\) by
\[
\kappa_6=169,
\qquad
\kappa_7=D_A,
\qquad
\kappa_{n+1}=D_A\kappa_n\quad(n\ge7).
\]
The first equality is the regional impulse; the next two describe the normalized continuation.

\begin{theorem}[Graduated analytic realization]
\label{p3-83-c30-s1:thm:realizacion-analitica-graduada}
The analytic germ compatible with the rules of
\cref{p3-83-c30-s1:def:lector-regional-determinantal} is unique and equals
\[
\boxed{
\mathscr C_\pi(z)
=169z^6+\frac{D_Az^7}{1-D_Az}.}
\]
Their coefficients satisfy
\[
\kappa_n=D_A^{n-6}\qquad(n\ge7).
\]
\end{theorem}

\begin{proof}
Write
\[
F(z)=169z^6+\sum_{k\ge1}c_kz^{6+k}.
\]
The normalization of continuation and the determinantal covariance give
\[
c_1=D_A,
\qquad
c_{k+1}=D_Ac_k.
\]
By induction, \(c_k=D_A^k\). The geometric sum yields
\[
\sum_{k\ge1}D_A^kz^{6+k}
=\frac{D_Az^7}{1-D_Az}.
\]
No free coefficients remain.
\end{proof}

Since \(0<\alpha<1\) and
\(0<D_A<1\),
\[
0<D_A\alpha<1.
\]
Consequently, \(\mathscr C_\pi(\alpha)\) converges absolutely and its
radius of convergence is \(D_A^{-1}>1\). Numerically,
\[
D_A=0.7751291192471564301888840964802\ldots,
\]
\[
D_A\alpha=0.00565639046986492673885079095469\ldots.
\]
The speed of this contraction explains why the complete sum and its
degree-seven truncation have the same first fifteen digits in the
final chart.

% HMT_SOURCE_SEGMENT_END id=B06
\subsection{Persistence of return in the
\texorpdfstring{\(5\)}{5}
regions and compatibility of the
\texorpdfstring{\(\eta_{\mathrm{ret}}\)}{eta_ret}
component}
\label{sec:s102:c38:2:5}

% HMT_SOURCE_SEGMENT_BEGIN id=A02 source=colaboracion/parte_iii/source/public_final/ascensos_20260826/16f_realizacion_analitica_contraangulo.tex body_lines=49-170

Let
\[
\mathcal F_\pi=\Psi^{-1}(010211)\subset\mathcal U_6
\]
the microfiber relative to the catalog. Recall its five elements, now
separated into initial block and terminal block:
\[
\begin{aligned}
 &(501,614,498\mid169,272,272),\\
 &(810,923,870\mid169,272,272),\\
 &(810,983,810\mid169,272,272),\\
 &(870,923,810\mid169,272,272),\\
 &(870,923,810\mid418,674,521).
\end{aligned}
\]
We define the return projection.
\[
p_{\rm ret}:\mathbb Z^6\longrightarrow\mathbb Z^3,
\qquad
p_{\rm ret}(u_1,\ldots,u_6)=(u_4,u_5,u_6).
\]

\begin{proposition}[Persistent terminal block]
\label{p3-20260826:prop:bloque-terminal-persistente}
The multiset \(p_{\rm ret}(\mathcal F_\pi)\) has a unique element of
maximum multiplicity:
\[
b_\pi=(169,272,272),
\qquad
\operatorname{mult}_{\mathcal F_\pi}(b_\pi)=4.
\]
The only remaining terminal block is
\[
b_\pi^-=(418,674,521)
\]
and has multiplicity one.
\end{proposition}

\begin{proof}
The assertion follows by exact enumeration of the five rows of the
catalog that satisfy \(\Psi(U)=010211\). Four rows have terminal
block \((169,272,272)\); the fifth, which carries the mutated return, has
\((418,674,521)\). The maximum multiplicity is four and is attained only
once.
\end{proof}

The oriented difference between the mutated return and the persistent is
\[
 \omega_\pi=b_\pi^- - b_\pi=(249,402,249),
 \qquad
 \langle(1,1,1),\omega_\pi\rangle=900.
\]
This identity separates two objects that must not be conflated. The integer
\(169\) is a coordinate of a terminal datum, that is, a datum of degree
zero. The vector \(\omega_\pi\) can be read as a one-cochain
after orienting the edge joining the two blocks. Calling it a cocycle would additionally
require fixing a cochain complex and verifying that its coboundary vanishes.
The analytic prolongation constructed below does not presuppose that
identification.

The selection of \(169\) does not require privileging the first position of the block. The symmetric group \(\mathfrak S_3\) acts by permuting its three coordinates. The stabilizer of \(b_\pi\) exchanges the two occurrences of \(272\) and fixes the unique coordinate of multiplicity one. Denote that coordinate by \(\operatorname{sing}(b)\) when the block has type \((r,s,s)\), with \(r\ne s\). Equivalently,
\[
\operatorname{sing}(b)
=\sum_{j=1}^3b_j-2\operatorname{moda}(b).
\]
In particular,
\[
\boxed{\operatorname{sing}(b_\pi)=169.}
\]

\begin{definition}[Residual fiber selector]
\label{p3-20260826:def:selector-residual-pi}
The selector \(\rho_\pi\) successively applies two operations:
\begin{enumerate}
  \item selects the terminal block of maximum multiplicity within
  \(p_{\rm ret}(\mathcal F_\pi)\);
  \item retains the coordinate fixed by the stabilizer of that block.
\end{enumerate}
By \cref{p3-20260826:prop:bloque-terminal-persistente},
\[
\boxed{\rho_\pi=169.}
\]
\end{definition}

\begin{theorem}[Equivariant uniqueness of return]
\label{p3-20260826:thm:unicidad-retorno-169}
Among the selectors that
\begin{enumerate}
  \item are invariant under permutations of the five regions;
  \item are equivariant under permutations of the three terminal
  coordinates;
  \item select the block of maximal multiplicity;
  \item retain a coordinate fixed by its stabilizer,
\end{enumerate}
the return value is uniquely determined and is \(169\).
\end{theorem}

\begin{proof}
Invariance under \(\mathfrak S_5\) prevents using the order of the rows. The
third condition selects \(b_\pi\), which is the unique mode by
\cref{p3-20260826:prop:bloque-terminal-persistente}. Its stabilizer in
\(\mathfrak S_3\) has two orbits on the coordinates: the pair of
values \(272\) and the singleton value \(169\). The fourth condition selects the
second orbit. Equivariance preserves its value when its position is shifted.
\end{proof}

This argument improves the positional reading of the return. The \(169\) is not
«the first coordinate» of a selected row: it is an invariant of the
regional multiset and of the action of its reordering symmetries.
Nor is it selected by global rarity. Among the \(468\) emissions in the
catalogue, \(169\) appears \(84\) times, uniformly distributed with
\(14\) occurrences in each of the six positions. The uniqueness of
\cref{p3-20260826:thm:unicidad-retorno-169} is therefore relational and regional: it resides
in the combination of microfiber, modal persistence, and stabilizer, not in the
isolated frequency of the integer.

% HMT_SOURCE_SEGMENT_END id=A02
\subsection{Application of the regional reader to the conjugate angular momentum pair}
\label{sec:s102:c38:2:6}

% HMT_SOURCE_SEGMENT_BEGIN id=A05 source=colaboracion/parte_iii/source/public_final/ascensos_20260826/16f_realizacion_analitica_contraangulo.tex body_lines=277-317

Analytic continuation is fixed by the following rules.

\begin{definition}[Regional Determinantal Reader]
\label{p3-20260826:def:lector-regional-determinantal}
The determinantal regional reader satisfies:
\begin{enumerate}
  \item \emph{degree compatibility}: cylindrical length is the analytic
  degree;
  \item \emph{renewal decomposition}: the finite return impulse and
  strict continuation belong to distinct additive summands;
  \item \emph{continuation normalization}: after recording the residue
  \(\rho_\pi\) just once, the surviving branch begins at the unit of
  the determinant line;
  \item \emph{concatenation covariance}: each elementary prolongation
  acts through \(\bigwedge^2\mathcal T=D_A\);
  \item \emph{orientation}: in the cardinal convention fixed for APP, the
  regional return belongs to the compensating sector and is subtracted from the fifth-degree
action phase.
\end{enumerate}
\end{definition}

The third rule has mathematical content. The value \(169\) is recorded once as the amplitude of the return; it is not multiplied again at each prolongation. The strict continuation counts the sole line that continues and normalizes it by its unit. This is the distinction between a renewal recurrence and homogeneous iteration of the initial impulse.

The need to declare that normalization can be proved. For every
\(\lambda\in\mathbb R\), the family
\[
F_\lambda(z)
=169z^6+\lambda\frac{D_Az^7}{1-D_Az}
\]
preserves the same finite return and the same determinantal ratio between
consecutive coefficients. The catalogue and the ratio alone do not distinguish
\(\lambda\). The unit of the determinant line fixes
\(\lambda=1\) before \(z=\alpha\) is evaluated. Thus, the normalization
is not obtained by fitting the final value; it forms part of the definition of the
character.

% HMT_SOURCE_SEGMENT_END id=A05
\subsection{Closed sum and publication of the posterior component upon return}
\label{sec:s102:c38:2:7}

% HMT_SOURCE_SEGMENT_BEGIN id=A06 source=colaboracion/parte_iii/source/public_final/ascensos_20260826/16f_realizacion_analitica_contraangulo.tex body_lines=319-380

We define the coefficients \(\kappa_n\) by
\[
\kappa_6=169,
\qquad
\kappa_7=D_A,
\qquad
\kappa_{n+1}=D_A\kappa_n\quad(n\ge7).
\]
The first equality is the regional impulse; the next two describe the normalized continuation.

\begin{theorem}[Graduated analytic realization]
\label{p3-20260826:thm:realizacion-analitica-graduada}
The analytic germ compatible with the rules of
\cref{p3-20260826:def:lector-regional-determinantal} is unique and equals
\[
\boxed{
\mathscr C_\pi(z)
=169z^6+\frac{D_Az^7}{1-D_Az}.}
\]
Their coefficients satisfy
\[
\kappa_n=D_A^{n-6}\qquad(n\ge7).
\]
\end{theorem}

\begin{proof}
Write
\[
F(z)=169z^6+\sum_{k\ge1}c_kz^{6+k}.
\]
The normalization of continuation and the determinantal covariance give
\[
c_1=D_A,
\qquad
c_{k+1}=D_Ac_k.
\]
By induction, \(c_k=D_A^k\). The geometric sum yields
\[
\sum_{k\ge1}D_A^kz^{6+k}
=\frac{D_Az^7}{1-D_Az}.
\]
No free coefficients remain.
\end{proof}

Since \(0<\alpha<1\) and
\(0<D_A<1\),
\[
0<D_A\alpha<1.
\]
Consequently, \(\mathscr C_\pi(\alpha)\) converges absolutely and its
radius of convergence is \(D_A^{-1}>1\). Numerically,
\[
D_A=0.7751291192471564301888840964802\ldots,
\]
\[
D_A\alpha=0.00565639046986492673885079095469\ldots.
\]
The speed of this contraction explains why the complete sum and its
degree-seven truncation have the same first fifteen digits in the
final chart.

% HMT_SOURCE_SEGMENT_END id=A06
\section{Confluence of \texorpdfstring{\(\eta_{\mathrm{ret}}\)}{eta_ret} With the valuation \texorpdfstring{\(\varepsilon_H=-34\)}{epsilon_H=-34} in the magnitude \texorpdfstring{\(\hbar_{\mathrm{ret}}\)}{hbar_ret}}
\label{sec:s102:c38:3}

\subsection{Derivation of the determinantal character of parangular transport}
\label{sec:s102:c38:3:1}

% HMT_SOURCE_SEGMENT_BEGIN id=B04 source=colaboracion/parte_iii/source/public_final/base_83/c30_body.tex body_lines=196-259

The dodecaphasic closure provides \(\alpha\). The completion of the three
nonadic pairs provides the scale \(10^3\), so that
\[
A=10^3\alpha.
\]
Having chosen the Archimedean angular chart, the common coordinate in radians is
\[
x=\frac{\pi A}{180}=\frac{50\pi}{9}\alpha.
\]
Let \(R\) be a real involution of zero trace in dimension two,
\[
R^2=I_2,
\qquad
\operatorname{tr}R=0,
\]
and consider the formal two-layer transport
\[
\mathcal T(x,y)=\exp(-xI_2-yR).
\]
The oriented variable \(y\) has not yet been evaluated. Nevertheless, the action
of \(\mathcal T\) on the determinant line is independent of it:
\[
\begin{aligned}
\det\mathcal T(x,y)
&=\exp\!\left(\operatorname{tr}(-xI_2-yR)\right)\\
&=e^{-2x}.
\end{aligned}
\]
We define
\[
\boxed{
D_A=e^{-2x}=e^{-100\pi\alpha/9}.}
\]
The subscript recalls that this is the determinant of the common contraction
associated with \(A\), not a function of \(C_*^{\rm HMT}\).

\begin{proposition}[Character of the determinant line]
\label{p3-83-c30-s1:prop:caracter-determinante}
The map
\[
\delta_A:(\mathbb N_0,+)\longrightarrow(\mathbb R_{>0},\cdot),
\qquad
\delta_A(k)=D_A^k,
\]
is the unique multiplicative character whose value on an elementary
prolongation is \(D_A\).
\end{proposition}

\begin{proof}
The multiplicativity gives
\[
\delta_A(k)
=\delta_A(\underbrace{1+\cdots+1}_{k\text{ times}})
=\delta_A(1)^k=D_A^k.
\]
The same identity proves existence and uniqueness.
\end{proof}

The determinant remains invariant under conjugation of
\(\mathcal T\) and under the exchange of sheets \(R\mapsto-R\). Therefore, the
tail constructed from \(D_A\) belongs to the common mode. The orientation
will appear in the sign with which the regional character corrects the phase.

% HMT_SOURCE_SEGMENT_END id=B04
\subsection{action character of degree \texorpdfstring{\(5\)}{5} in the construction regional}
\label{sec:s102:c38:3:2}

% HMT_SOURCE_SEGMENT_BEGIN id=B07 source=colaboracion/parte_iii/source/public_final/base_83/c30_body.tex body_lines=366-405

The central block
\[
\mathcal B_c=\begin{pmatrix}7&2\\2&7\end{pmatrix}
\]
has eigenvalues \(9\) and \(5\). The dodecaphasic cycle contributes rank
\(12\); its step-three orbits have length four; and the complete census
contributes the exact quotient
\[
\frac{104976}{1944}=54.
\]
From these invariants the fifth-degree action character is fixed.
\[
\boxed{
H_5(\alpha,\varphi)
=\frac12\sqrt{\frac{1000\alpha}{\varphi}}-\alpha
+\frac9{16}\alpha^2-\frac59\alpha^3
+\frac7{48}\alpha^4-\frac1{54}\alpha^5.}
\]
The four rational coefficients can be written as
\[
\frac9{16}=\frac{\lambda_+}{4^2},
\qquad
\frac59=\frac{\lambda_-}{\lambda_+},
\qquad
\frac7{48}=\frac{\tfrac12\operatorname{tr}\mathcal B_c}{4\cdot12},
\qquad
\frac1{54}=\frac{1944}{104976},
\]
with \((\lambda_+,\lambda_-)=(9,5)\). These identities certify the
provenance of the numbers used. The successive assignment of those
invariants to degrees two, three, four, and five forms part of the analytic
character; it does not follow from the mere existence of the finite catalog.

This precision removes a common ambiguity. A combination of
invariants of the construction is an internal and reproducible definition; its
naturality requires the reader rules specifying order, signs, and
normalization. Here those rules are visible and are subject to the same
negative controls as the rest of the construction.

% HMT_SOURCE_SEGMENT_END id=B07
\subsection{Exact decomposition of the action line into pre-return and post-return sections}
\label{sec:s102:c38:3:3}

% HMT_SOURCE_SEGMENT_BEGIN id=B12 source=colaboracion/parte_iii/source/public_final/base_83/c30_body.tex body_lines=571-617

The fifth-degree character and the post-action component after return
are two distinct values:
\[
H_5
=1.05457181764615446962582182943306\ldots,
\]
\[
\bar h_{\rm HMT}^{\rm ret}
=1.05457181691503906113369058472430\ldots.
\]
The first is the local action before incorporating the regional character; the
second is the oriented action remaining after the return. The correction
that generates the conjugate angular coordinate prevents identifying them.

The metrological comparison must respect this distinction. The first value gives
the mantissa
\[
2\pi H_5
=6.62607014999998792600111034989947\ldots,
\]
while the subsequent action gives
\[
2\pi\bar h_{\rm HMT}^{\rm ret}
=6.62607014540625433351074984759288\ldots.
\]
In the International System, the value
\[
h=6.62607015\times10^{-34}\;\mathrm{J\,s}
\]
is exact by definition \cite{BIPMSI}. By therefore,
\[
2\pi H_5-6.62607015
=-1.2073998889650101\ldots\times10^{-14},
\]
and it is not an exact equality. The proximity of the value before the return is
a remarkable numerical check; it does not authorize replacing the defining constant
of the SI with the subsequent component or confusing the two HMT actions\@.

This result precisely determines the mathematical situation. The regional
reader generates \(C_*^{\rm HMT}\) and, from it, the functional
\(\Gamma_{6\to8}\). The same calculation separates the Planck approximation
produced by \(H_5\) from the reduced action intervening in
the conjugate angular coordinate. The separation does not weaken the construction: it eliminates an
identification incompatible with the exact figures and converts every
comparison into a falsifiable assertion of a well-defined type.

% HMT_SOURCE_SEGMENT_END id=B12
\subsection{Factorization of the determinantal character in the \texorpdfstring{\(2\)}{2} phase charts}
\label{sec:s102:c38:3:4}

% HMT_SOURCE_SEGMENT_BEGIN id=A04 source=colaboracion/parte_iii/source/public_final/ascensos_20260826/16f_realizacion_analitica_contraangulo.tex body_lines=211-275

Dodecaphase closure supplies \(\alpha\). The declared angular insertion
uses the scale \(10^3\), compatible with the grouping of three nonadic
pairs, so that
\[
\Adeg=10^3\alpha\,{}^\circ.
\]
Having chosen the Archimedean angular chart, the common coordinate in radians is
\[
x=: \Arad=\frac{\pi\Adeg}{180}=\frac{50\pi}{9}\alpha.
\]
Let \(R\) be a real involution of zero trace in dimension two,
\[
R^2=I_2,
\qquad
\operatorname{tr}R=0,
\]
and consider the formal two-layer transport
\[
\mathcal T(x,y)=\exp(-xI_2-yR).
\]
The oriented variable \(y\) has not yet been evaluated. Nevertheless, the action
of \(\mathcal T\) on the determinant line is independent of it:
\[
\begin{aligned}
\det\mathcal T(x,y)
&=\exp\!\left(\operatorname{tr}(-xI_2-yR)\right)\\
&=e^{-2x}.
\end{aligned}
\]
We define
\[
\boxed{
D_A=e^{-2x}=e^{-100\pi\alpha/9}.}
\]
The subscript recalls that this is the determinant of the common contraction
associated with \(A\), not a function of the counter-angle.

\begin{proposition}[Character of the determinant line]
\label{p3-20260826:prop:caracter-determinante}
The map
\[
\delta_A:(\mathbb N_0,+)\longrightarrow(\mathbb R_{>0},\cdot),
\qquad
\delta_A(k)=D_A^k,
\]
is the unique multiplicative character whose value on an elementary
prolongation is \(D_A\).
\end{proposition}

\begin{proof}
The multiplicativity gives
\[
\delta_A(k)
=\delta_A(\underbrace{1+\cdots+1}_{k\text{ times}})
=\delta_A(1)^k=D_A^k.
\]
The same identity proves existence and uniqueness.
\end{proof}

The determinant remains invariant under conjugation of
\(\mathcal T\) and under the exchange of sheets \(R\mapsto-R\). Therefore, the
tail constructed from \(D_A\) belongs to the common mode. The orientation
will appear in the sign with which the regional character corrects the phase.

% HMT_SOURCE_SEGMENT_END id=A04
\subsection{Publication of the action character of degree \texorpdfstring{\(5\)}{5} in the parangular chart}
\label{sec:s102:c38:3:5}

% HMT_SOURCE_SEGMENT_BEGIN id=A07 source=colaboracion/parte_iii/source/public_final/ascensos_20260826/16f_realizacion_analitica_contraangulo.tex body_lines=382-418

The central block
\[
\mathcal B_c=\begin{pmatrix}7&2\\2&7\end{pmatrix}
\]
has eigenvalues \(9\) and \(5\). The dodecaphasic cycle contributes rank
\(12\); its step-three orbits have length four; and the complete census
contributes the exact quotient
\[
\frac{104976}{1944}=54.
\]
The current reader declares the nature of fifth-degree action
\[
\boxed{
H_5(\alpha,\varphi)
=\frac12\sqrt{\frac{1000\alpha}{\varphi}}-\alpha
+\frac9{16}\alpha^2-\frac59\alpha^3
+\frac7{48}\alpha^4-\frac1{54}\alpha^5.}
\]
The four rational coefficients can be written as
\[
\frac9{16}=\frac{\lambda_+}{4^2},
\qquad
\frac59=\frac{\lambda_-}{\lambda_+},
\qquad
\frac7{48}=\frac{\tfrac12\operatorname{tr}\mathcal B_c}{4\cdot12},
\qquad
\frac1{54}=\frac{1944}{104976},
\]
with \((\lambda_+,\lambda_-)=(9,5)\). These identities fix the
internal arithmetic representability of the character coefficients. Their
order, signs, and normalization belong to the graded rule of the regional
reader and are applied before any chart of units. The ablations
certify that every term effectively intervenes in the resulting phase.
In particular, neither the SI mantissa of action nor the sexagesimal value of the
counter-angle belongs to the domain of the generating function.

% HMT_SOURCE_SEGMENT_END id=A07
\section{Parangular application \texorpdfstring{\(P(\alpha,\eta_{\mathrm{ret}})\)}{P(alpha,eta_ret)} and bilayer closure \texorpdfstring{\(C^*=2(\eta_{\mathrm{ret}}+\alpha)\)}{C*=2(eta_ret+alpha)}}
\label{sec:s102:c38:4}

\subsection{Oriented phase and conjugate angular coordinate}
\label{sec:s102:c38:4:1}

% HMT_SOURCE_SEGMENT_BEGIN id=B08 source=colaboracion/parte_iii/source/public_final/base_83/c30_body.tex body_lines=407-472

The fifth-degree phase prior to return is
\[
y_5=\frac\pi{90}\bigl(H_5+\alpha\bigr).
\]
The complete regional character produces the correction
\(\mathscr C_\pi(\alpha)\). With the orientation fixed in
\cref{p3-83-c30-s1:def:lector-regional-determinantal}, we define
\[
\boxed{
y_*=\frac\pi{90}\bigl(H_5+\alpha\bigr)
-169\alpha^6
-\frac{D_A\alpha^7}{1-D_A\alpha}.}
\]
This is the primary formula. The sexagesimal unit does not come into play until applying the chart.
\[
C_*^{\rm HMT}=\frac{180}{\pi}y_*.
\]

\begin{theorem}[Conjugate angular coordinate of the regional reader determinantal]
\label{p3-83-c30-s1:thm:contraangulo-regional}
With the dodecaphasic closure of \(\alpha\),
\(\varphi=(1+\sqrt5)/2\) and the rules of the determinantal regional reader fixed, the
phase \(y_*\) and its conjugate angular coordinate are determined without introducing the sought
value. Their evaluation is
\[
\boxed{
y_*=0.03706622646583826398493779409230638\ldots,}
\]
\[
\boxed{
C_*^{\rm HMT}
=2.12373833896864572426195138039336\ldots{}^\circ.}
\]
Therefore, to fifteen decimal places,
\[
\boxed{C_*^{\rm HMT}=2.123738338968646^\circ.}
\]
\end{theorem}

\begin{proof}
\cref{p3-83-c30-s1:thm:unicidad-retorno-169} determines \(169\). The closure of
\(\alpha\) determines \(A\), \(x\), and
\(D_A=e^{-100\pi\alpha/9}\). \cref{p3-83-c30-s1:thm:realizacion-analitica-graduada}
determines the complete series. The formula for \(H_5\) contains only
\(\alpha\), \(\varphi\), and declared structural invariants. Substituting
these data into the expression for \(y_*\) gives the printed value. The decimal value
of \(C_*^{\rm HMT}\) occurs in none of the right-hand sides.
\end{proof}

The reduced action component after return is now as follows
output:
\[
\boxed{
\bar h_{\rm HMT}^{\rm ret}
=\frac{C_*^{\rm HMT}}2-\alpha
=H_5-\frac{90}{\pi}\mathscr C_\pi(\alpha).}
\]
Numerically,
\[
\bar h_{\rm HMT}^{\rm ret}
=1.05457181691503906113369058472430\ldots.
\]
The inverse identity has not been used to introduce \(C_*^{\rm HMT}\); it is
obtained after generating \(y_*\).

% HMT_SOURCE_SEGMENT_END id=B08
\subsection{Conjugate coordinate \texorpdfstring{\(C^*\)}{C*} and reversal of its orientation under the bilayer involution}
\label{sec:s102:c38:4:2}

% HMT_SOURCE_SEGMENT_BEGIN id=A08 source=colaboracion/parte_iii/source/public_final/ascensos_20260826/16f_realizacion_analitica_contraangulo.tex body_lines=420-503

The fifth-degree phase prior to return is
\[
y_5=\frac\pi{90}\bigl(H_5+\alpha\bigr).
\]
The complete regional character produces the correction
\(\mathscr C_\pi(\alpha)\). With the orientation fixed in
\cref{p3-20260826:def:lector-regional-determinantal}, we define
\[
\boxed{
\Crad:=y_*=\frac\pi{90}\bigl(H_5+\alpha\bigr)
-169\alpha^6
-\frac{D_A\alpha^7}{1-D_A\alpha}.}
\]
From now on we also write
\[
 \boxed{y_{\rm reg}:=y_*}
\]
when it is necessary to distinguish this realization from the spectral phase of
\Gtr. The two notations here designate the same canonical regional object.
This is the primary formula. The sexagesimal unit does not intervene until the chart is applied
\[
\boxed{\Cdeg:=\frac{180}{\pi}\Crad.}
\]
The symbols \(\Crad\) and \(\Cdeg\) do not denote two counterangles: they are the
radial and sexagesimal coordinates of the same angular element.

\begin{theorem}[Counter-angle of the regional determinantal reader]
\label{p3-20260826:thm:contraangulo-regional}
Given the dodecaphase closure of \(\alpha\),
\(\varphi=(1+\sqrt5)/2\), the germ \(H_5\) and the rules of the regional
determinantal reader, the phase \(\Crad\) and its counter-angle are determined without
introducing the sought value into the evaluation function. Its evaluation is
\[
\boxed{
\Crad=0.03706622646583826398493779409230638\ldots,}
\]
\[
\boxed{
\Cdeg
=2.12373833896864572426195138039336\ldots{}^\circ.}
\]
Therefore, to fifteen decimal places,
\[
\boxed{\Cdeg=2.123738338968646^\circ.}
\]
\end{theorem}

\begin{proof}
The \cref{p3-20260826:thm:unicidad-retorno-169} determines \(169\). The closure of \(\alpha\) determines \(\Adeg\), \(\Arad\), and \(D_A=e^{-100\pi\alpha/9}\). The \cref{p3-20260826:thm:realizacion-analitica-graduada} determines the complete series. The fixed formula for \(H_5\) contains only \(\alpha\), \(\varphi\), and declared structural invariants. Substituting these data into the expression for \(\Crad\) yields the printed value. The decimal value of the counter-angle appears in none of the right-hand sides of this evaluation. Comparison with a metrological action mantissa is undertaken only after the phase has been obtained and does not alter the generating function.
\end{proof}

Let \(u_\circ\) be the sexagesimal base whose complete turn has coordinate
\(360\), and let us introduce the angular coordinate
\(\alpha_{\angle,{\rm deg}}:=\alpha_{\rm HMT}\). The output posterior to
return is then typed by
\[
\boxed{
\etaRet
=\frac{\Cdeg}2-\alpha_{\angle,{\rm deg}}
=H_5-\frac{90}{\pi}\mathscr C_\pi(\alpha).}
\]
Its radial coordinate is
\[
 \etaRetrad=\frac{\pi}{180}\etaRet
 =\frac{\Crad}{2}-\frac{\pi}{180}\alpha_{\rm HMT}.
\]
Numerically,
\[
\etaRet
=1.05457181691503906113369058472430\ldots.
\]
The inverse identity has not been used to introduce \(\Cdeg\); it is obtained
after generating \(\Crad\). The symbol \(\hbar\) is reserved for the
element of the action line constructed in the following section.

% HMT_SOURCE_SEGMENT_END id=A08
\subsection{Direct theorem \texorpdfstring{\(C^*=2(\eta_{\mathrm{ret}}+\alpha)\)}{C*=2(eta_ret+alpha)} and inverse character of the identity \texorpdfstring{\(\eta_{\mathrm{ret}}=C^*/2-\alpha\)}{eta_ret=C*/2-alpha}}
\label{sec:s102:c38:4:3}

% HMT_SOURCE_SEGMENT_BEGIN id=A13 source=colaboracion/parte_iii/source/public_final/ascensos_20260826/16f_realizacion_analitica_contraangulo.tex body_lines=750-797

The fifth-degree character and the coordinate after return
are two distinct values:
\[
H_5
=1.05457181764615446962582182943306\ldots,
\]
\[
\etaRet
=1.05457181691503906113369058472430\ldots.
\]
The former is the local character before incorporation of the regional character; the
latter is the oriented coordinate remaining after the return. The correction
generating the contra-angle prevents their identification.

The metrological comparison must respect this distinction. The first value gives
the mantissa
\[
2\pi H_5
=6.62607014999998792600111034989947\ldots,
\]
while the posterior coordinate gives
\[
2\pi\etaRet
=6.62607014540625433351074984759288\ldots.
\]
As an external control only, we retain the mantissa \(6.62607015\) of the
defining value of the International System \cite{BIPMSI}. Neither its exponent nor its unit is yet incorporated into the HMT line: the decade is derived in the
\cref{chap:orden-determinantal-accion} and the unit belongs to a later
metrological chart. The comparison of mantissas gives
\[
2\pi H_5-6.62607015
=-1.2073998889650101\ldots\times10^{-14},
\]
and is not an exact equality. The proximity of the pre-return value is
a noteworthy numerical control; it does not authorize replacing the defining SI
constant by the subsequent coordinate or confusing an angular coordinate with
the action element that will be constructed upon incorporating the decade.

This result determines the mathematical situation precisely. The regional
reader generates \((\Crad,\Cdeg)\) and, from that single angular element, the functional
\(\Gamma_{6\to8}\). The same calculation separates the external proximity of the
mantissa produced by \(H_5\) from the return coordinate that enters the
counter-angle. The separation does not weaken the construction: it eliminates an
identification incompatible with the exact digits and turns each
comparison into a falsifiable assertion of a well-defined type.

% HMT_SOURCE_SEGMENT_END id=A13
% HMT_SOURCE_SEGMENT_BEGIN id=D01 source=manuscrito/sucesor_102/deltas_ley9/c38_par_angular_cartas_y_canales.tex body_lines=3-20
\label{sec:l9s102:par-angular}

The cubic completion
\begin{equation}
 10^3=(9+1)^3=9^3+3\cdot9^2+3\cdot9+1
 =729+243+27+1
 \label{eq:l9s102:base-1000}
\end{equation}
determines the chart global of the coordinate direct. The closure bilayer contributes
the factor \(2\) of the coordinate conjugate:
\begin{equation}
 \boxed{
 A_{\deg}=10^3\alpha_{\mathrm{HMT}},
 \qquad
 C^*_{\deg}=2(\eta_{\mathrm{ret}}+\alpha_{\mathrm{HMT}})
 \quad\text{in }\mathbb R/(360\mathbb Z).}
 \label{eq:l9s102:a-cstar-deg}
\end{equation}
% HMT_SOURCE_SEGMENT_END id=D01
% HMT_SOURCE_SEGMENT_BEGIN id=D04 source=manuscrito/sucesor_102/deltas_ley9/c38_par_angular_cartas_y_canales.tex body_lines=60-66
The inverse identity
\begin{equation}
 \eta_{\mathrm{ret}}=\frac12C^*_{\deg}-\alpha_{\mathrm{HMT}}
 \label{eq:l9s102:eta-corolario-inverso}
\end{equation}
is established as an algebraic corollary of the direct construction.

% HMT_SOURCE_SEGMENT_END id=D04
\section{Exact commutativity of the charts in turns, degrees, and radians}
\label{sec:s102:c38:5}

\subsection{Derivation of the cylindrical length and the analytic degree of the parangular precursor}
\label{sec:s102:c38:5:1}

% HMT_SOURCE_SEGMENT_BEGIN id=B03 source=colaboracion/parte_iii/source/public_final/base_83/c30_body.tex body_lines=157-194

Let \(\mathcal P\) be the free module generated by the regional words and
let \(F^n\mathcal P\) be its length filtration. The associated grading
records the first level at which each word appears. For a parameter
\(z\), the ordinary length character is
\[
\chi_z([w])=z^{|w|}.
\]
Multiplicativity
\[
\chi_z([uv])=\chi_z([u])\chi_z([v])
\]
expresses that concatenation adds lengths. When evaluated at
\(z=\alpha\), a word of length \(n\) receives degree \(\alpha^n\).

\begin{definition}[Degree compatibility]
\label{p3-83-c30-s1:def:compatibilidad-grado}
The regional analytic reader is compatible with the cylindrical filtration
when it sends the homogeneous component of length \(n\) to analytic degree
\(n\):
\[
\operatorname{gr}_n\mathcal P\longrightarrow
\alpha^n\mathbb R.
\]
\end{definition}

The five regions belong to \(\mathcal U_6\); therefore, the return datum
appears in degree six:
\[
\rho_\pi[U_6]\longmapsto169\alpha^6.
\]
This six is the length of the regional word. It is not the length of the
closed walks of \cref{chap:cinco-ciclos-pi}. Those walks traverse
the four anchors
\(U_\pi,U_e,U_\varphi,U_{\rm APP}\) simultaneously and have minimum length eight. A
word, an edge of the block graph, and a closed walk are objects of
distinct types; the grading used here belongs to the first.

% HMT_SOURCE_SEGMENT_END id=B03
\subsection{Publication of the analytic degree on the global phase chart}
\label{sec:s102:c38:5:2}

% HMT_SOURCE_SEGMENT_BEGIN id=A03 source=colaboracion/parte_iii/source/public_final/ascensos_20260826/16f_realizacion_analitica_contraangulo.tex body_lines=172-209

Let \(\mathcal P\) be the free module generated by the regional words and
let \(F^n\mathcal P\) be its length filtration. The associated grading
records the first level at which each word appears. For a parameter
\(z\), the ordinary length character is
\[
\chi_z([w])=z^{|w|}.
\]
Multiplicativity
\[
\chi_z([uv])=\chi_z([u])\chi_z([v])
\]
expresses that concatenation adds lengths. When evaluated at
\(z=\alpha\), a word of length \(n\) receives degree \(\alpha^n\).

\begin{definition}[Degree compatibility]
\label{p3-20260826:def:compatibilidad-grado}
The regional analytic reader is compatible with the cylindrical filtration
when it sends the homogeneous component of length \(n\) to analytic degree
\(n\):
\[
\operatorname{gr}_n\mathcal P\longrightarrow
\alpha^n\mathbb R.
\]
\end{definition}

The five regions belong to \(\mathcal U_6\); therefore, the return datum
appears in degree six:
\[
\rho_\pi[U_6]\longmapsto169\alpha^6.
\]
This six is the length of the regional word. It is not the length of the
closed walks of \cref{chap:cinco-ciclos-pi}. Those walks traverse
the four anchors
\(U_\pi,U_e,U_\varphi,U_{\rm APP}\) simultaneously and have minimum length eight. A
word, an edge of the block graph, and a closed walk are objects of
distinct types; the grading used here belongs to the first.

% HMT_SOURCE_SEGMENT_END id=A03
% HMT_SOURCE_SEGMENT_BEGIN id=D02 source=manuscrito/sucesor_102/deltas_ley9/c38_par_angular_cartas_y_canales.tex body_lines=21-34
The global calendar exhibits the three factorizations
\begin{equation}
 360=12\cdot30=4\cdot90=3\cdot120.
 \label{eq:l9s102:360-factorizaciones}
\end{equation}
The analytical chart is obtained through
\begin{equation}
 A_{\mathrm{rad}}=\frac\pi{180}A_{\deg}
 =\frac{50\pi}{9}\alpha_{\mathrm{HMT}},
 \qquad
 C^*_{\mathrm{rad}}=\frac\pi{180}C^*_{\deg}
 =\frac\pi{90}(\eta_{\mathrm{ret}}+\alpha_{\mathrm{HMT}}).
 \label{eq:l9s102:a-cstar-rad}
\end{equation}
% HMT_SOURCE_SEGMENT_END id=D02
% HMT_SOURCE_SEGMENT_BEGIN id=D03 source=manuscrito/sucesor_102/deltas_ley9/c38_par_angular_cartas_y_canales.tex body_lines=35-59

\begin{theorem}[Commutativity of Regional and Action Factorizations]
\label{thm:l9s102:conmutatividad-cstar}
The two routes toward the conjugate phase produce the same output:
\begin{equation}
 \boxed{
 \frac\pi{180}C^*_{\deg}
 =\frac\pi{90}(\eta_{\mathrm{ret}}+\alpha_{\mathrm{HMT}})
 =\frac\pi{90}(H_5+\alpha_{\mathrm{HMT}})
 -\mathscr C_\pi(\alpha_{\mathrm{HMT}})=y_*.}
 \label{eq:l9s102:cuadrado-cstar}
\end{equation}
\end{theorem}

\begin{proof}
Substituting \cref{eq:l9s102:eta-ret} in the second member yields
\[
 \frac\pi{90}
 \left(H_5-\frac{90}{\pi}\mathscr C_\pi+\alpha_{\mathrm{HMT}}\right)
 =\frac\pi{90}(H_5+\alpha_{\mathrm{HMT}})-\mathscr C_\pi.
\]
The final expression is the regional definition of \(y_*\), whereas the
first equality follows from \cref{eq:l9s102:a-cstar-deg}.
\end{proof}

% HMT_SOURCE_SEGMENT_END id=D03
\section{spectral decomposition \texorpdfstring{\(q_\pm\)}{q+/-} and calculo functional on the channels \texorpdfstring{\(90/120\)}{90/120} previously generated}
\label{sec:s102:c38:6}

\subsection{Spectral reconstruction of the angular pair and hexad-octad transition}
\label{sec:s102:c38:6:1}

% HMT_SOURCE_SEGMENT_BEGIN id=B11 source=colaboracion/parte_iii/source/public_final/base_83/c30_body.tex body_lines=523-569

The common phase and the oriented phase determine the two channels.
\[
q_-=e^{-x+y_*},
\qquad
q_+=e^{-x-y_*}.
\]
The invariants are recovered exactly.
\[
q_-q_+=D_A,
\qquad
\alpha=-\frac9{100\pi}\log D_A,
\qquad
C_*^{\rm HMT}=\frac{90}{\pi}\log\frac{q_-}{q_+}.
\]
The first equality reads the common contraction; the third, the oriented
separation of the sheets. The exchange \(R\mapsto-R\) permutes
\(q_+\) and \(q_-\), preserves \(D_A\) and reverses the sign of \(y_*\).

For
\[
S_{90}(q)=\frac{q^{90}}{1-q^{270}},
\qquad
S_{120}(q)=\frac{q^{120}}{1-q^{360}},
\]
the hexad--octad transition functional is written entirely in the
phase chart:
\[
\Gamma_{6\to8}
=\frac{180}{\pi}xy_*
+12\bigl(S_{90}(q_-)-S_{90}(q_+)\bigr)
-\bigl(S_{120}(q_-)-S_{120}(q_+)\bigr).
\]
The evaluation produced by the coordinate \(C_*^{\rm HMT}\) of
\cref{p3-83-c30-s1:thm:contraangulo-regional} is
\[
\boxed{
\Gamma_{6\to8}
=0.274007672694459274747255260774\ldots.}
\]
Thus, the internal value recognized in Dimensional Mechanics as the Barbero--Immirzi functional ceases to receive \(C_*^{\rm HMT}\) as an independent literal: both descend from the same pair of phases generated in this chapter. The geometric and physical interpretation of \(\Gamma_{6\to8}\) belongs to the work on Dimensional Mechanics; its mathematical evaluation already belongs to the previously constructed interface.

% HMT_SOURCE_SEGMENT_END id=B11
\subsection{Relative transporter and spectral disjunction of the blocks}
\label{sec:s102:c38:6:2}

% HMT_SOURCE_SEGMENT_BEGIN id=B13 source=colaboracion/parte_iii/source/public_final/base_83/c30_body.tex body_lines=619-675
\label{p3-83-c30-s1:sec:transportador-relativo-split}

In the ordered basis of channels, let
\[
 R=\begin{pmatrix}0&1\\1&0\end{pmatrix},
 \qquad P_\pm=\frac{I\pm R}{2},
 \qquad
 B_0=7I+2R,
 \qquad B_1=\Adeg I+\Cdeg R.
\]
The split decomposition
\[
 xI+yR=\rho(x,y)e^{\eta(x,y)R},
 \quad \rho(x,y)=\sqrt{x^2-y^2},
 \quad \eta(x,y)=\operatorname{artanh}(y/x)
\]
constructs the unique transporter
\[
 T_{\rm rel}
 =\frac{\sqrt{(\Adeg)^2-(\Cdeg)^2}}{\sqrt{45}}
   \exp\!\left[
   \left(\operatorname{artanh}\frac{\Cdeg}{\Adeg}
   -\operatorname{artanh}\frac27\right)R\right],
 \qquad T_{\rm rel}B_0=B_1.
\]
In the spectral basis it acts by
\(9\mapsto\Adeg+\Cdeg\) and
\(5\mapsto\Adeg-\Cdeg\).

\begin{theorem}[Spectral Disjunction of the Local and Angular Realizations]
\label{p3-83-c30-s1:thm:no-go-entrelazador-bloques-split}
The spectra
\[
 \operatorname{Spec}(B_0)=\{9,5\},\qquad
 \operatorname{Spec}(B_1)=\{\Adeg+\Cdeg,\Adeg-\Cdeg\}
\]
are disjoint. Consequently,
\[
 \operatorname{Hom}_{B_0,B_1}
 :=\{\Phi:\Phi B_0=B_1\Phi\}=\{0\}.
\]
\end{theorem}

\begin{proof}
The rational bounds constructed for the angular coordinates give
\(7.29<\Adeg<7.30\) and \(2.12<\Cdeg<2.13\); therefore,
\(9.41<\Adeg+\Cdeg<9.43\) and
\(5.16<\Adeg-\Cdeg<5.18\). In the basis that diagonalizes \(R\), each entry
of an intertwiner is multiplied by the difference between an eigenvalue
of \(B_0\) and one of \(B_1\); disjointness forces all of them to vanish.
\end{proof}

The relative transporter compares two blocks within the split algebra;
the dodecaphasic--Witt and incidence intertwiners act between
equivalent representations in their own spaces of residence. This separation
positively types both operations and preserves their domains.

% HMT_SOURCE_SEGMENT_END id=B13
\subsection{Spectral evaluation \texorpdfstring{\(V_{90,120}\)}{V_90,120} of the incidence channels previously generated by the TPK}
\label{sec:s102:c38:6:3}

% HMT_SOURCE_SEGMENT_BEGIN id=A11 source=colaboracion/parte_iii/source/public_final/ascensos_20260826/16f_realizacion_analitica_contraangulo.tex body_lines=554-572

The common phase and the oriented phase determine the two channels.
\[
q_-=e^{-\Arad+\Crad},
\qquad
q_+=e^{-\Arad-\Crad}.
\]
The invariants are recovered exactly.
\[
q_-q_+=D_A,
\qquad
\alpha=-\frac9{100\pi}\log D_A,
\qquad
\Cdeg=\frac{90}{\pi}\log\frac{q_-}{q_+}.
\]
The first equality reads the common contraction; the third, the oriented
separation of the sheets. The exchange \(R\mapsto-R\) permutes
\(q_+\) and \(q_-\), preserves \(D_A\) and reverses the sign of \(\Crad\).

% HMT_SOURCE_SEGMENT_END id=A11
% HMT_SOURCE_SEGMENT_BEGIN id=D05 source=manuscrito/sucesor_102/deltas_ley9/c38_par_angular_cartas_y_canales.tex body_lines=68-97
\label{sec:l9s102:qpm-posteriores}

Let \(H^2=I\) and \(P_\pm=(I\pm H)/2\). The angular block
\begin{equation}
 B=A_{\mathrm{rad}}I+C^*_{\mathrm{rad}}H
 =(A_{\mathrm{rad}}+C^*_{\mathrm{rad}})P_+
 +(A_{\mathrm{rad}}-C^*_{\mathrm{rad}})P_-
 \label{eq:l9s102:b-angular}
\end{equation}
is published multiplicatively by
\begin{equation}
 e^{-B}=q_+P_++q_-P_-,
 \qquad
 q_\pm=\exp[-(A_{\mathrm{rad}}\pm C^*_{\mathrm{rad}})].
 \label{eq:l9s102:qpm-angular}
\end{equation}
The involution \(C^*\mapsto-C^*\) permutes \(q_+\) and \(q_-\), and preserves
\begin{equation}
 q_+q_-=e^{-2A_{\mathrm{rad}}}=D_A.
 \label{eq:l9s102:det-angular}
\end{equation}
The cardinalities \(90/120\), previously constructed through TPK
incidence, now receive their spectral evaluation:
\begin{equation}
 V_{90,120}
 =(I-T^{90})(I-T^{120})^{-1}
 =\sum_{\sigma\in\{+,-\}}
 \frac{1-q_\sigma^{90}}{1-q_\sigma^{120}}P_\sigma.
 \label{eq:l9s102:v90120-posterior}
\end{equation}
% HMT_SOURCE_SEGMENT_END id=D05
\subsubsection{Split relative transporter}
\label{sec:s102:c38:6:3:1}

% HMT_SOURCE_SEGMENT_BEGIN id=A12 source=colaboracion/parte_iii/source/public_final/ascensos_20260826/16f_realizacion_analitica_contraangulo.tex body_lines=574-748
\label{p3-20260826:sec:transportador-relativo-split}

The comparison With the central block of APP can be carried out without identifying
representations distinct. Fijemos the ordered basis of channels and the
involution
\[
 R=\begin{pmatrix}0&1\\1&0\end{pmatrix},
 \qquad R^2=I,
 \qquad
 P_\pm=\frac{I\pm R}{2}.
\]
In commutative algebra
\[
 \mathscr B_R:=\{xI+yR:x,y\in\mathbb R\}
\]
consider a, within the already declared sexagesimal chart,
\[
 B_0:=7I+2R,
 \qquad
 B_1:=\Adeg I+\Cdeg R.
\]
The orthogonal change of basis
\[
 S=\frac1{\sqrt2}
 \begin{pmatrix}1&1\\1&-1\end{pmatrix}
\]
simultaneously diagonalizes the involution and both blocks. For \(x>|y|\),
we define
\[
 \rho(x,y):=\sqrt{x^2-y^2},
 \qquad
 \eta(x,y):=\operatorname{artanh}\!\left(\frac yx\right).
\]
Then
\[
 xI+yR=\rho(x,y)e^{\eta(x,y)R}.
\]
If
\[
 \rho_0=\sqrt{45},\quad
 \eta_0=\operatorname{artanh}(2/7),\quad
 \rho_1=\rho(\Adeg,\Cdeg),\quad
 \eta_1=\eta(\Adeg,\Cdeg),
\]
the unique left multiplication that carries \(B_0\) to \(B_1\) is
\[
 \boxed{
 T_{\rm rel}
 :=\frac{\rho_1}{\sqrt{45}}
 e^{(\eta_1-\eta_0)R},
 \qquad
 T_{\rm rel}B_0=B_1.}
\]

\begin{proposition}[Relative carrier of the two blocks]
\label{p3-20260826:prop:transportador-relativo-split}
Once \(R\), the channel order, the positive branch, and the sexagesimal
chart have been fixed, \(T_{\rm rel}\) is the unique element of
\(\mathscr B_R\) satisfying \(TB_0=B_1\). In the spectral basis it acts by
\[
 9\longmapsto\Adeg+\Cdeg,
 \qquad
 5\longmapsto\Adeg-\Cdeg.
\]
\end{proposition}

\begin{proof}
The split polar decomposition yields the printed formula and proves existence.
Since \(\det B_0=45\ne0\), every matrix solution of \(TB_0=B_1\) satisfies
\(T=B_1B_0^{-1}\); it is therefore unique. Diagonalization by means of \(S\)
produces the two spectral quotients
\((\Adeg+\Cdeg)/9\) and \((\Adeg-\Cdeg)/5\), equivalent to the exponential
expression of \(T_{\rm rel}\).
\end{proof}

This transporter is not a representational intertwiner. Indeed, an
intertwiner \(\Phi\) between the endomorphisms \(B_0\) and \(B_1\) should
satisfy
\[
 \Phi B_0=B_1\Phi.
\]

\begin{theorem}[Spectral Obstruction to Entanglement]
\label{p3-20260826:thm:no-go-entrelazador-bloques-split}
For the values of \(\Adeg\) and \(\Cdeg\) constructed in this section,
\[
 \boxed{\Phi B_0=B_1\Phi\quad\Longrightarrow\quad\Phi=0.}
\]
\end{theorem}

\begin{proof}
The spectra are
\[
 \operatorname{Spec}(B_0)=\{9,5\},
 \qquad
 \operatorname{Spec}(B_1)
 =\{\Adeg+\Cdeg,\Adeg-\Cdeg\}.
\]
The previously demonstrated rational bounds provide the disjoint intervals.
\[
 7.29<\Adeg<7.30,
 \qquad
 2.12<\Cdeg<2.13,
\]
and, consequently,
\[
 9.41<\Adeg+\Cdeg<9.43,
 \qquad
 5.16<\Adeg-\Cdeg<5.18.
\]
In the basis that diagonalizes \(R\), each entry \(\phi_{ij}\) of \(\Phi\)
is multiplied by the difference between an eigenvalue of \(B_0\) and one of
\(B_1\). All these differences are nonzero; hence all entries of
\(\Phi\) vanish.
\end{proof}

Numerically,
\[
 \rho_1=6.9814819335172378048\ldots,
 \qquad
 \frac{\rho_1}{\sqrt{45}}=1.0407378791354140779\ldots,
\]
\[
 \eta_1-\eta_0
 =0.005796351470571295590\ldots.
\]
The formula preserves the split algebra, the projectors \(P_\pm\), and the
Lorentzian form up to the conformal factor \(\rho_1^2/45\). It nevertheless depends
on \(B_1\) already constructed: it does not constitute a prior derivation of
\(\Adeg\) or \(\Cdeg\), nor should it be confused with the
equivariant intertwiners between realizations of the same module.

For
\[
S_{90}(q)=\frac{q^{90}}{1-q^{270}},
\qquad
S_{120}(q)=\frac{q^{120}}{1-q^{360}},
\]
the hexad--octad transition functional is written entirely in the
phase chart:
\[
\Gamma_{6\to8}
=\frac{180}{\pi}\Arad\Crad
+12\bigl(S_{90}(q_-)-S_{90}(q_+)\bigr)
-\bigl(S_{120}(q_-)-S_{120}(q_+)\bigr).
\]
The evaluation produced by the counter-angle of the
\cref{p3-20260826:thm:contraangulo-regional} is
\[
\boxed{
\Gamma_{6\to8}
=0.274007672694459274747255260774\ldots.}
\]
The mathematically equivalent forms of the internal functional are
\[
\boxed{
\Gamma_{6\to8}
=\Arad\Cdeg+K_{6\to8}
=\frac{180}{\pi}\Arad\Crad+K_{6\to8}
=\frac{\pi\Adeg\Cdeg}{180}+K_{6\to8}.}
\]
These identities evaluate \(\Gamma_{6\to8}\) after generating
\((\Crad,\Cdeg)\); it is not used in the inverse sense to select the
counter-angle. The physical reading in the Ashtekar--Barbero connection is
then represented by means of
\[
 \Gamma_{6\to8}\dashrightarrow\gamma_{\rm BI}.
\]
The arrow introduces no third functional and does not act retroactively on \(C^*\). The
series \(S_{90}\), \(S_{120}\), the coefficients \(12\) and \(-1\), and their
assignment to the \(6\to8\) passage belong to the declared mathematical
construction. Physical identification additionally requires the connection's own
normalization.

% HMT_SOURCE_SEGMENT_END id=A12
\section{Naturalness under Refinement and Separation with Respect to Metrological Recognition}
\label{sec:s102:c38:7}

\subsection{Exact estimation of the remainder of the analytic prolongation}
\label{sec:s102:c38:7:1}

% HMT_SOURCE_SEGMENT_BEGIN id=B09 source=colaboracion/parte_iii/source/public_final/base_83/c30_body.tex body_lines=474-499

The expression initially found,
\[
y_7=\frac\pi{90}(H_5+\alpha)-169\alpha^6-D_A\alpha^7,
\]
is the seventh-degree truncation of the full character. Produce
\[
C_7=2.12373833896864600265403827103919\ldots{}^\circ.
\]
The omitted remainder is not an indeterminate term; it is the exact geometric tail.
\[
\mathscr C_\pi(\alpha)
-169\alpha^6-D_A\alpha^7
=\frac{D_A^2\alpha^8}{1-D_A\alpha}.
\]
Therefore,
\[
\left|C_7-C_*^{\rm HMT}\right|
=\frac{180}{\pi}\frac{D_A^2\alpha^8}{1-D_A\alpha}
=2.7839208689064583\ldots\times10^{-16}\;{}^\circ.
\]
The recurrence thus transforms a seventh-degree approximation into a
complete analytic realization and provides a closed-form bound for each
subsequent truncation.

% HMT_SOURCE_SEGMENT_END id=B09
\subsection{Functional dependence of the regional selector and the analytic continuation}
\label{sec:s102:c38:7:2}

% HMT_SOURCE_SEGMENT_BEGIN id=B10 source=colaboracion/parte_iii/source/public_final/base_83/c30_body.tex body_lines=501-521

The construction changes when any of its active components is removed.
The following values are obtained without recalibrating the other terms:
\[
\begin{array}{@{}lc@{}}
\toprule
\text{variant}&\text{conjugate angular coordinate in degrees}\\
\midrule
\rho_\pi=168&2.1237383389772976863904487\ldots\\
\rho_\pi=169&2.1237383389686457242619514\ldots\\
\rho_\pi=170&2.1237383389599937621334541\ldots\\
\text{without determinantal tail}&2.1237383389686949415301675\ldots\\
D_A\text{ replaced by }1&2.1237383389686313409965826\ldots\\
\bottomrule
\end{array}
\]
These ablations do not in themselves constitute a probabilistic estimate.
They demonstrate that the equivariant selector and the determinantal character are
effective components of the output and cannot be removed while maintaining the
same realization.

% HMT_SOURCE_SEGMENT_END id=B10
\subsection{Operatorial Interpretation of the Conjugate Angular Coordinate}
\label{sec:s102:c38:7:3}

% HMT_SOURCE_SEGMENT_BEGIN id=B14 source=colaboracion/parte_iii/source/public_final/base_83/c30_body.tex body_lines=677-715

The entire procedure can be summarized as a chain of applications:
\[
\begin{aligned}
\mathcal F_\pi
&\longmapsto b_\pi
\longmapsto\rho_\pi=169,\\
\alpha
&\longmapsto A
\longmapsto x
\longmapsto D_A,\\
(\rho_\pi,D_A)
&\longmapsto\mathscr C_\pi(\alpha),\\
(H_5,\mathscr C_\pi(\alpha))
&\longmapsto y_*
\longmapsto C_*^{\rm HMT},\\
(x,y_*)
&\longmapsto(q_+,q_-)
\longmapsto\Gamma_{6\to8}.
\end{aligned}
\]
The first line belongs to the finite catalog; the second, to the common closure of
\(\alpha\); the third, to the determinantal prolongation; the fourth produces the
oriented coordinate; the fifth transports it toward the dimensional interface.

In this organization, \(\alpha\) measures the common contraction of the transport, and \(C_*^{\rm HMT}\) measures the oriented separation of its two sheets. The regional return adds no external constant: it selects a finite amplitude and prolongs it through the determinant already fixed by \(\alpha\). The chart in degrees is subsequent. The primary object is the pair of phases \((x,y_*)\).

An exhaustive verification reconstructs the five regions directly
from the catalogue, traverses the actions of
\(\mathfrak S_5\) and \(\mathfrak S_3\), checks the recurrence, and sums the
series. The input of \(\alpha\) is the exact output of its
dodecaphasic closure, not a transcribed metrological digit. The verification keeps
the catalogue theorems, reader rules, and external comparison
separate; therefore, no experimental comparison intervenes in the generating
function.
% HMT_SOURCE_SEGMENT_END id=B14
\subsection{Angular Truncation Convergence Certification}
\label{sec:s102:c38:7:4}

% HMT_SOURCE_SEGMENT_BEGIN id=A09 source=colaboracion/parte_iii/source/public_final/ascensos_20260826/16f_realizacion_analitica_contraangulo.tex body_lines=505-530

The expression initially found,
\[
y_7=\frac\pi{90}(H_5+\alpha)-169\alpha^6-D_A\alpha^7,
\]
is the seventh-degree truncation of the full character. Produce
\[
C_7=2.12373833896864600265403827103919\ldots{}^\circ.
\]
The omitted remainder is not an indeterminate term; it is the exact geometric tail.
\[
\mathscr C_\pi(\alpha)
-169\alpha^6-D_A\alpha^7
=\frac{D_A^2\alpha^8}{1-D_A\alpha}.
\]
Therefore,
\[
\left|C_7-\Cdeg\right|
=\frac{180}{\pi}\frac{D_A^2\alpha^8}{1-D_A\alpha}
=2.7839208689064583\ldots\times10^{-16}\;{}^\circ.
\]
The recurrence thus transforms a seventh-degree approximation into a
complete analytic realization and provides a closed-form bound for each
subsequent truncation.

% HMT_SOURCE_SEGMENT_END id=A09
\subsection{Naturality of the regional selector under refinement and chart change}
\label{sec:s102:c38:7:5}

% HMT_SOURCE_SEGMENT_BEGIN id=A10 source=colaboracion/parte_iii/source/public_final/ascensos_20260826/16f_realizacion_analitica_contraangulo.tex body_lines=532-552

The construction changes when any of its active components is removed.
The following values are obtained without recalibrating the other terms:
\[
\begin{array}{@{}lc@{}}
\toprule
\text{variant}&\text{counter-angle in degrees}\\
\midrule
\rho_\pi=168&2.1237383389772976863904487\ldots\\
\rho_\pi=169&2.1237383389686457242619514\ldots\\
\rho_\pi=170&2.1237383389599937621334541\ldots\\
\text{without determinantal tail}&2.1237383389686949415301675\ldots\\
D_A\text{ replaced by }1&2.1237383389686313409965826\ldots\\
\bottomrule
\end{array}
\]
These ablations do not in themselves constitute a probabilistic estimate.
They demonstrate that the equivariant selector and the determinantal character are
effective components of the output and cannot be removed while maintaining the
same realization.

% HMT_SOURCE_SEGMENT_END id=A10
\subsection{Metric independence of the parangular construction}
\label{sec:s102:c38:7:6}

% HMT_SOURCE_SEGMENT_BEGIN id=A14 source=colaboracion/parte_iii/source/public_final/ascensos_20260826/16f_realizacion_analitica_contraangulo.tex body_lines=799-831

The entire procedure can be summarized as a chain of applications:
\[
\begin{aligned}
\mathcal F_\pi
&\longmapsto b_\pi
\longmapsto\rho_\pi=169,\\
\alpha
&\longmapsto \Adeg
\longmapsto \Arad
\longmapsto D_A,\\
(\rho_\pi,D_A)
&\longmapsto\mathscr C_\pi(\alpha),\\
(H_5,\mathscr C_\pi(\alpha))
&\longmapsto \Crad
\longmapsto \Cdeg,\\
(\Arad,\Crad)
&\longmapsto(q_+,q_-)
\longmapsto\Gamma_{6\to8}.
\end{aligned}
\]
The first line belongs to the finite catalog; the second, to the common closure of
\(\alpha\); the third, to the determinantal prolongation; the fourth produces the
oriented coordinate; the fifth transports it toward the dimensional interface.

In this organization, \(\alpha\) measures the common contraction of transport, and
the counter-angle measures the oriented separation of its two sheets. In the
current evaluation, the regional return selects a finite amplitude and
prolongs it through the determinant already fixed by \(\alpha\), without inserting an
additional external constant. The degree chart is subsequent. The primary
object is the phase pair \((\Arad,\Crad)\). The recurrence, regional sum,
and ablations depend on the catalogue and on
\(\alpha_{\rm HMT}\), but receive no additional metrological constant.
% HMT_SOURCE_SEGMENT_END id=A14

